%% ma: L12-B03-0010
%% lop: 12
%% chuong: 1
%% bai: 3
%% dang: 12.3.8
%% loai: TN4
%% muc: TH
%% dap_an: "B"
%% nguon: "Ảnh giáo viên gửi 11/10/2026 (Đề kiểm tra giữa HK1 Toán 12, THPT Việt Đức 2026-2027, mã đề 0103), Phần I Câu 11"
%% kien_thuc: "Bài 3 (định nghĩa tiệm cận đứng, tiệm cận ngang qua giới hạn)"
%% trang_thai: cho-duyet
%% ly_do: "Dạng mới đề xuất, chờ giáo viên duyệt."
\begin{ex}
Cho hàm số $y=f(x)$ xác định trên tập $\mathbb R\setminus\{1\}$ và liên tục trên từng khoảng xác định. Biết rằng $\lim\limits_{x\to+\infty}f(x)=2$; $\lim\limits_{x\to-\infty}f(x)=-2$; $\lim\limits_{x\to1^+}f(x)=+\infty$. Tổng số đường tiệm cận đứng và ngang của đồ thị hàm số $y=f(x)$ là
\choice{$1$}{\True $3$}{$4$}{$2$}
\loigiai{$\lim\limits_{x\to+\infty}f(x)=2$ và $\lim\limits_{x\to-\infty}f(x)=-2$ nên đồ thị có hai tiệm cận ngang $y=2$ và $y=-2$. $\lim\limits_{x\to1^+}f(x)=+\infty$ nên đường thẳng $x=1$ là tiệm cận đứng. Tổng cộng có $3$ đường tiệm cận.}
\end{ex}
